% called by main.tex
%
\chapter{Despliegue con Docker en local}
\label{anexo::docker-local}

% [PENDIENTE: verificar los comandos exactos contra README.md y docs/GUIA_USO.md del
%  repositorio antes de la entrega; añadir requisitos de versión de Docker si procede.]

Este anexo describe cómo construir y ejecutar el sistema en una máquina local a partir
del código fuente, mediante Docker. La contenedorización garantiza un entorno aislado y
reproducible, evitando los problemas de compatibilidad de dependencias.

\section{Requisitos previos}
\label{anexo::docker-local_req}

Es necesario tener instalado \textbf{Docker} (y \textit{Docker Compose}, incluido en
Docker Desktop). Conviene disponer del \textit{checkpoint} del modelo en la carpeta
\texttt{models/} del repositorio; sin él, la API arranca con un error descriptivo.

\section{Construcción y ejecución}
\label{anexo::docker-local_run}

Una vez clonado el repositorio, desde su raíz se construyen y levantan los dos
servicios (API e interfaz web) con una única orden:

\begin{verbatim}
docker-compose up --build
\end{verbatim}

La imagen se construye sobre \texttt{python:3.12-slim} e instala las versiones de
PyTorch \textbf{para CPU}. Tras el arranque, los servicios quedan accesibles en:

\begin{itemize}
  \item \textbf{API REST}: \url{http://localhost:8000} (documentación interactiva
        OpenAPI en \texttt{/docs}).
  \item \textbf{Interfaz web}: \url{http://localhost:8501}.
\end{itemize}

Los modelos y la configuración se montan como volúmenes, de modo que no es necesario
reconstruir la imagen al cambiar de \textit{checkpoint}. Para detener los servicios:

\begin{verbatim}
docker-compose down
\end{verbatim}

% [PENDIENTE: incluir aquí, si se quiere, la lista completa de comandos de la guía
%  (entrenamiento, evaluación, pre-resize, bootstrap) o remitir al anexo/capítulo
%  correspondiente.]
